\section{Why an Id Has to Say What It Names}
\label{sec:ch04-why-an-id-has-to-say-what-it-names}

\index{Relay!object identification specification|(}%
The fix has a specification behind it, and the specification is two paragraphs
long. \emph{GraphQL Global Object Identification} asks a server for exactly two
things~\autocite{relayobjectid}. First, an interface:

\begin{quote}
The server must provide an interface called \code{Node}. That interface must
include exactly one field, called \code{id} that returns a non-null \code{ID}.
\end{quote}

And second, a way in:

\begin{quote}
The server must provide a root field called \code{node} that returns the
\code{Node} interface. This root field must take exactly one argument, a
non-null ID named \code{id}.
\end{quote}

\index{node@\code{node}!what it buys a client}%
Read what that gives a client rather than what it asks of you. Any object it
holds, of any type, it can refetch by handing back the string the server gave
it, without knowing which field it originally came through or which service
answered. A cache can refresh one entry; a page can deep-link to a session
whose type it worked out at run time. The alternative is a lookup field per
type, which is \code{sessionById} multiplied by every type in the graph.

\index{ID@\code{ID}!is not an integer}%
\code{ID} is a scalar in its own right and it is not \code{Int} wearing a
different name. It serializes as a string, it accepts a string, and the
specification says nothing whatever about what is inside it. The words
\enquote{opaque} and \enquote{base64} appear nowhere in that document. What it
requires is that the value be globally unique and that handing it back
refetches the same object.

\index{ID@\code{ID}!opacity is a convention}%
Opacity is a convention that grew up around the requirement rather than part
of it, and it is worth keeping for a reason that is not about tidiness. An id
a client can take apart is an id a client will take apart, and the moment one
does, the encoding is a public API you did not mean to publish. Base64 is weak
enough that anyone determined will read it in a second, and awkward enough
that nobody does it by accident.

\index{ID@\code{ID}!is not a secret}%
So an opaque id is not a secret. Base64 is an encoding, not a cipher, and
anything you put in one you have published: an internal customer number, or
the fact that your ids are sequential and therefore how many sessions you
have. Chapter~\ref{ch:entities-authorization} is about a related mistake one
level up, where a route into the graph is assumed to be guarded because the
field in front of it is.
\index{Relay!object identification specification|)}%
